\SourceLeaf{247}%
the magnitude of the line element on a curved surface. Let us first examine the geometric meaning of the coefficients \(E,F,G\). We already noted in Article 5 that we can conceive of two systems of lines on a curved surface. Along each line of the first system only \(p\) varies and \(q\) is constant. Along each line of the second system only \(q\) varies and \(p\) is constant. Any point of the surface can be regarded as the intersection of a line of the first system with a line of the second. The element of the first line adjacent to this point and corresponding to the variation \(\mathrm{d}p\) is \(=\BareRoot{E}\mathord{.}\mathrm{d}p\), and the element of the second line corresponding to the variation \(\mathrm{d}q\) is \(=\BareRoot{G}\mathord{.}\mathrm{d}q\). Finally, if we denote the angle between these elements by \(\omega\), it is easy to see that \(\cos\omega=\frac{F}{\BareRoot{EG}}\). The area of the parallelogram element on the curved surface between the two lines of the first system corresponding to \(q,q+\mathrm{d}q\) and the two lines of the second system corresponding to \(p,p+\mathrm{d}p\) is \(\BareRoot{(EG-FF)}\mathrm{d}p\mathord{.}\mathrm{d}q\).

Any line on the curved surface belonging to neither system arises when we regard \(p\) and \(q\) as functions of a single new variable or one of them as a function of the other. Let \(s\) be the length of such a curve, measured from an arbitrary starting point, with either direction taken as positive. Let \(\theta\) denote the angle that the element \(\mathrm{d}s=\BareRoot{(E\mathrm{d}p^2+2F\mathrm{d}p\mathord{.}\mathrm{d}q+G\mathrm{d}q^2)}\) makes with the line of the first system drawn through the start of the element. To leave no ambiguity, we assume that the angle is always measured from the branch of this line along which \(p\) increases, and that it is taken as positive toward the side on which \(q\) increases. With these conventions it is easy to see that
\[
\begin{aligned}
\cos\theta\mathord{.}\mathrm{d}s&=\BareRoot{E}\mathord{.}\mathrm{d}p+\BareRoot{G}\mathord{.}\cos\omega\mathord{.}\mathrm{d}q=\frac{E\mathrm{d}p+F\mathrm{d}q}{\BareRoot{E}}\\
\sin\theta\mathord{.}\mathrm{d}s&=\BareRoot{G}\mathord{.}\sin\omega\mathord{.}\mathrm{d}q=\frac{\BareRoot{(EG-FF)}\mathord{.}\mathrm{d}q}{\BareRoot{E}}
\end{aligned}
\]
\Article{18}
We shall now investigate the condition under which this line is shortest. Since its length \(s\) is given by the integral
\[
s=\int\BareRoot{(E\mathrm{d}p^2+2F\mathrm{d}p\mathord{.}\mathrm{d}q+G\mathrm{d}q^2)}
\]
the condition for a minimum requires that the variation of this integral arising from an infinitesimal change in the path of the line be \(=0\). For our purpose the calculation is easier in this case if we regard \(p\) as a function of \(q\).
